% ============================================================
% Draft section for "Luxury good taxation proposals"
% Slots in after Section 7.2 (or replaces it)
% ============================================================

\section{Implementation: an ad valorem excise above category-specific per-item price thresholds}
\label{sec:implementation}

\subsection{Why an excise rather than an additional VAT rate}

Section~\ref{sec:reform} presented the EU-wide additive luxury tax as a new VAT rate. Two features
of the VAT Directive make an excise the more natural legal form, notwithstanding the administrative
advantages of VAT.

First, the Directive contains no legal basis for an \emph{increased} rate of any kind. Articles
96--98 provide for a standard rate of at least 15 per cent and up to two reduced rates drawn from
Annex III; the increased rates that several member states operated before the single market ---
France applied a 33\textfrac{1}{3} per cent luxury rate to car sales until September 1987, and
Italy ran rates above its standard rate into the 1980s --- were abolished in the rate-approximation
exercise of the early 1990s. Introducing an EU-wide luxury VAT is therefore not a matter of adding
a category to Annex III but of reinstating a rate tier that was deliberately removed, which is an
amendment to the rate architecture itself.

Second, a VAT rate is by construction a single percentage applied to the taxable amount. Any
liability that depends on the level of the price --- a threshold, a bracket schedule, or the
Danish-style progressive registration tax --- cannot be expressed as a VAT rate, and must take the
form of an excise. Since a Council directive is required in either case, the choice between the two
legal forms is not a choice about the required political effort. It is a choice about design
freedom, and the excise form buys a great deal of it at low cost.

The administrative-cost argument of Section~\ref{sec:reform} survives this change. An excise
assessed by VAT-registered vendors on the first supply to a non-registered person, computed on the
VAT-exclusive taxable amount and remitted on the VAT return, inherits the registration, invoicing
and audit infrastructure of the VAT system. What it does not inherit is the input-credit chain, and
it does not need to: the tax is designed to fall once, at the point of final supply or importation.

\subsection{The proposed structure}

Let $\mathcal{C}$ denote a set of luxury categories, each defined by an enumerated list of CN
eight-digit codes.\footnote{The Combined Nomenclature (Council Regulation (EEC) No 2658/87,
republished annually) is the EU's own eight-digit goods classification: the first six digits are the
WCO Harmonised System heading and subheading, the seventh and eighth an EU subdivision. Defining an
excise base by CN code is established practice --- Directive 92/83/EEC defines beer, wine and
spirits by CN heading, and the Energy Taxation Directive lists taxable energy products the same
way. Greece already defines its luxury tax base this way (Section~\ref{sec:precedents}).} For each
category $c$ let $\bar{p}_c$ be a per-item price threshold, expressed net of VAT and of the luxury
excise itself, and let $t^{\mathrm{EU}}$ be a single EU-wide ad valorem rate applying uniformly
across all categories. For an item with pre-tax price $p$ the liability is

\begin{equation}
T_c(p) =
\begin{cases}
0, & p \leq \bar{p}_c,\\[2mm]
\min\left\{\, t^{\mathrm{EU}}\, p, \;\; \dfrac{m}{m-1}\, t^{\mathrm{EU}} \left(p - \bar{p}_c\right)\right\}, & p > \bar{p}_c,
\end{cases}
\label{eq:liability}
\end{equation}

with a phase-in parameter $m > 1$. The two branches coincide at $p = m\bar{p}_c$, so the schedule
has three regions: no liability below the threshold; a phase-in band $(\bar{p}_c,\, m\bar{p}_c)$ in
which the tax is levied at the rate $\tfrac{m}{m-1}t^{\mathrm{EU}}$ on the excess over the
threshold; and a uniform ad valorem rate $t^{\mathrm{EU}}$ on the full price for every item priced
at or above $m\bar{p}_c$. The average rate rises continuously from zero at $\bar{p}_c$ to
$t^{\mathrm{EU}}$ at $m\bar{p}_c$ and is flat thereafter. Setting $m=2$ and $t^{\mathrm{EU}}=0.2$
gives $T_c(p)=\min\{0.2\,p,\; 0.4\,(p-\bar{p}_c)\}$.

The essential property of \eqref{eq:liability} is that $T_c$ is continuous at $\bar{p}_c$. The
schedule is a kink, not a notch. This is the single most important design choice in this section,
and we motivate it in Section~\ref{sec:nonotch}. Everything else --- the choice of $t^{\mathrm{EU}}$,
the category list, the level of the thresholds --- can be varied without disturbing it.

Relative to a hypothetical schedule that applied $t^{\mathrm{EU}}$ to the full price of every item
above the threshold, the phase-in forgoes, per item in the band,
\begin{equation}
t^{\mathrm{EU}}\,\frac{m\bar{p}_c - p}{m-1},
\label{eq:forgone}
\end{equation}
which equals $t^{\mathrm{EU}}\bar{p}_c$ at the threshold and declines linearly to zero at
$m\bar{p}_c$. The aggregate cost is therefore governed by the mass of the price distribution in a
band of width $(m-1)\bar{p}_c$ immediately above the threshold, and can be made small either by
setting $m$ close to one or by placing $\bar{p}_c$ where the density is thin. It is not a large
number for plausible parameters, and it buys the elimination of the distortions discussed next.

\subsection{Why not a clean threshold}
\label{sec:nonotch}

A binary threshold --- the tax applying to the whole price once the item crosses $\bar{p}_c$ ---
creates a discontinuity in liability of size $t^{\mathrm{EU}}\bar{p}_c$. Three consequences follow.

\emph{Price bunching.} Applying the standard small-notch approximation, the marginal buncher's
response satisfies $\Delta p / \bar{p}_c \approx \sqrt{2\,e\,t^{\mathrm{EU}}}$, where $e$ is the
elasticity of within-category quality spending with respect to the net-of-tax price. At
$t^{\mathrm{EU}}=0.2$ this implies a frictionless bunching window of 45 per cent of the threshold
price for $e=0.5$ and 63 per cent for $e=1$. The relevant elasticity here is emphatically not the
category-level own-price elasticity reported in Table~\ref{tab:elasticities}: that table excludes
inter-brand substitution by construction, because inter-brand substitution is irrelevant to a
category-wide tax. A price threshold, by contrast, is a tax on \emph{which} item within the category
is bought, and inter-brand and inter-model substitution is precisely the margin it operates on. The
elasticities governing bunching are the larger ones.

\emph{Missing product varieties.} Because demand falls discretely as the price crosses the
threshold, no producer will position a product in a band immediately above it. The distortion is
not merely a relocation of prices but a hole in the product space, and its welfare cost falls on the
consumers whose preferred variety is no longer offered. The empirical literature on attribute-based
notches makes the magnitude concrete. Ito and Sallee (2018) find that Japan's weight-based
fuel-economy steps induced around 10 per cent of vehicles to increase weight, by an average of
110 kg, and that the resulting distortion makes the second-best fuel-economy subsidy 40 per cent
lower than the first-best. India's 4{,}000 mm length threshold in vehicle excise --- roughly 29 per
cent below the line against 45 per cent above it --- did not merely shift prices but called into
existence an entire vehicle segment of sub-four-metre sedans and SUVs designed around the line.
Where price notches themselves have been studied, in real-estate transfer taxes, Kopczuk and Munroe
(2015) find that the volume of missing transactions above the notch exceeds the bunching mass below
it, and that incidence local to the threshold can exceed the value of the tax; Best and Kleven
(2018) find that eliminating a 1 per cent transaction tax raised market activity by 20 per cent.

\emph{Loss of signalling precision.} There is in addition a cost specific to the framework of
Section~\ref{sec:signalling}. If luxury consumption transmits information about income or wealth,
compressing a band of the price distribution onto a single focal point degrades the informativeness
of that signal in exactly the region where the tax binds, prompting substitution towards other
signalling margins. A notch thus loses corrective bite while retaining the distortion. The uniform
ad valorem rate of \eqref{eq:liability} above $m\bar{p}_c$ has the opposite property: it leaves the
ordering and relative spacing of observable prices intact and simply reduces the real resources
absorbed in transmitting the same information.

\subsection{Precedents}
\label{sec:precedents}

Taxes of this general shape are common, and two of the closest analogues are inside the EU.

\emph{Greece} levies a luxury tax (Law 3833/2010, Art.~17) on passenger cars of CN heading 8703 and
jeep-type vehicles of CN 8704, at ad valorem rates graduated by the manufacturer's wholesale price:
10 per cent from \euro{}17{,}000 to \euro{}22{,}000, 20 per cent above \euro{}22{,}000 to
\euro{}30{,}000, and 30 per cent above \euro{}30{,}000, assessed on the taxable value defined for
excise purposes. A separate Greek luxury tax applies to goods enumerated directly by CN code,
including furskin apparel (CN 4303), leather articles of exotic skins (CN 4202\,21\,00 and
4202\,31\,00), and silk-containing carpets. Greece thus already operates both elements we propose:
a CN-defined luxury base, and ad valorem rates keyed to a per-item pre-VAT price threshold. It is
also a cautionary example, since its brackets are notches rather than kinks.

\emph{France} levies the \emph{taxe forfaitaire sur les objets pr\'ecieux} (CGI Arts.~150\,VI--VJ)
at 6 per cent plus 0.5 per cent CRDS on sales of jewellery, watches, art, collectibles and
antiques, and 11 per cent plus CRDS on precious metals. Sales of jewellery, art, collectibles and
antiques are exempt where the price does not exceed \euro{}5{,}000. Two features are directly
relevant. The threshold is assessed \emph{objet par objet}, item by item, which is the rule we
propose; and it is assessed on the set rather than the item where the objects form a coherent
ensemble, which is the anti-fragmentation rule discussed in Section~\ref{sec:item}.

\emph{Canada}'s Select Luxury Items Tax Act, in force since 1 September 2022, is the closest
existing analogue to \eqref{eq:liability}. The tax is calculated as the lesser of 10 per cent of the
full price and 20 per cent of the amount above a threshold of CAD 100{,}000 for vehicles and
aircraft and CAD 250{,}000 for vessels --- that is, exactly \eqref{eq:liability} with $m=2$ and
$t^{\mathrm{EU}}=0.10$. It is collected from registered vendors on the first retail sale or on
importation. Its subsequent history is instructive in both directions: the levies on aircraft and
vessels were repealed with effect from 5 November 2025 on the grounds that they were inefficient and
costly to administer, while the levy on vehicles was retained and remains in force. The
vehicles-versus-vessels asymmetry is not surprising --- vessels are mobile, easily flagged and
berthed abroad, and thinly traded --- and it reinforces the case made in Section~\ref{sec:selfinterest}
for acting at EU rather than national level.

\emph{The United States} operated a luxury excise from 1991 under the Omnibus Budget Reconciliation
Act of 1990, at 10 per cent of the amount by which the price exceeded \$30{,}000 for cars,
\$100{,}000 for boats, \$250{,}000 for aircraft and \$10{,}000 for jewellery and furs. This was a
pure kink and generated no notch, and a GATT panel in \emph{United States --- Taxes on Automobiles}
(DS31/R, 1994, unadopted) found the price-threshold design consistent with Article~III:2. The
boats, aircraft, jewellery and furs components were repealed in 1993 and the automobile component
phased out by 2003. Two administrative lessons carried over into later drafting: the threshold was
indexed to inflation from 1994, and the statute had to be amended to exempt the subsequent
installation of small-value parts and accessories, a fragmentation response we return to below.

\emph{China} applies a 10 per cent surcharge to super-luxury vehicles priced above RMB 1.3 million,
levied at the retail stage on top of its category-based consumption tax --- a per-item price
threshold layered on a CN-style category list, as proposed here.

\subsection{Defining the item, and anti-fragmentation}
\label{sec:item}

A per-item threshold invites the splitting of a supply into components each priced below
$\bar{p}_c$. Three rules contain this at acceptable cost. First, the taxable item should be defined
by reference to the CJEU's single-supply doctrine (\emph{Card Protection Plan}, C-349/96;
\emph{Levob}, C-41/04), under which economically indivisible elements supplied together constitute
one supply. Second, following the French \emph{ensemble} rule, items supplied simultaneously to the
same customer that form a coherent set should be aggregated for the purpose of applying
$\bar{p}_c$. Third, following the US experience, accessories and parts installed within a defined
period after the first supply should be added to the taxable amount, with a de minimis exclusion to
avoid burdening ordinary after-sales servicing.

\subsection{Two parameters requiring explicit decision}

\emph{Interaction with the VAT base.} The revenue expressions of Section~\ref{sec:revenue} assume
that the EU-wide rate is additive, so that the total wedge is $1 + t_i + t^{\mathrm{EU}}$. If the
excise instead forms part of the taxable amount for VAT, the wedge compounds to
$(1+t_i)(1+t^{\mathrm{EU}})$. At the average statutory VAT rate of 22 per cent, an additive rate of
20 per cent and a compounded rate of 16.4 per cent both deliver the target total rate of 42 per
cent. The choice is immaterial to the economics but must be made explicitly, since the two produce
different headline rates for the same policy.

\emph{Services.} The per-item threshold is well defined for the goods categories of
Table~\ref{tab:base} --- automotive, personal luxury goods, fine wines and spirits, design and
furniture --- and for yachts. It is not well defined for hospitality (\euro{}54bn) and gourmet food
and fine dining (\euro{}18bn), which together account for roughly 22 per cent of the base. These
require either a threshold expressed per unit of service (per room-night, per cover) or a
category-level rule with no threshold. We adopt the former, with thresholds set so that the
resulting coverage is comparable to that in the goods categories.